\subsection{Absolute value of a complex measure}

Let \(\mu\) be a complex measure on a locally compact space \(X\) ; for every function \(f \in \mathcal{K}_{+}(X)\) , the positive real number

\[
L (f) = \sup _ {| g | \leqslant f, g \in \mathcal {K} (\mathrm{X}; \mathbf {C})} \left| \int g d \mu \right|\tag{11}
\]

is finite, because the relation \(|g| \leqslant f\) implies that \(\operatorname{Supp}(g) \subset \operatorname{Supp}(f)\) and \(\| g \| \leqslant \| f \|\) , thus our assertion follows from formula (4) of No.~3. Let us show that \(L\) can be extended, in only one way, to a positive measure on \(X\) ; in view of No.~5, Th. 1 and Ch. II, §2, No.~1, Prop. 3, it will suffice to show that if \(f_1, f_2\) are two functions in \(\mathcal{K}_+(X)\) , then \(L(f_1 + f_2) = L(f_1) + L(f_2)\) . Now, if \(|g_1| \leqslant f_1\) and \(|g_2| \leqslant f_2\) , where \(g_1\) and \(g_2\) are functions in \(\mathcal{K}(X; \mathbf{C})\) , we have \(|g_1 + \zeta g_2| \leqslant f_1 + f_2\) for any complex number \(\zeta\) of absolute value 1, therefore

\[
\left| \mu (g _ {1} + \zeta g _ {2}) \right| = \left| \mu (g _ {1}) + \zeta \mu (g _ {2}) \right| \leqslant L (f _ {1} + f _ {2}).
\]

Moreover, we can suppose \(\zeta\) so chosen that

\[
\left| \mu (g _ {1}) + \zeta \mu (g _ {2}) \right| = \left| \mu (g _ {1}) \right| + \left| \mu (g _ {2}) \right|;
\]

since \(|\mu(g_i)|\) is arbitrarily close to \(L(f_i)\) ( \(i = 1,2\) ), this proves that \(L(f_1) + L(f_2) \leqslant L(f_1 + f_2)\) . On the other hand, consider a function \(g \in \mathcal{K}(\mathbf{X};\mathbf{C})\) such that \(|g| \leqslant f_1 + f_2\) . The function \(g_i\) equal to \(g f_i / (f_1 + f_2)\) at the points where \(f_1(x) + f_2(x) \neq 0\) , and to 0 elsewhere ( \(i = 1,2\) ), is continuous on \(\mathbf{X}\) because \(f_i / (f_1 + f_2)\) ( \(i = 1,2\) ) is continuous at every point where \(f_1(x) + f_2(x) \neq 0\) and we have \(|g_i(x)| \leqslant |g(x)|\) for every \(x \in \mathbf{X}\) , which proves the continuity of \(g_i\) at the points where \(f_1(x) + f_2(x) = 0\) ( \(i = 1,2\) ), since at these points we have also \(g(x) = 0\) . It is clear that \(|g_i| \leqslant f_i\) ( \(i = 1,2\) ) and \(g = g_1 + g_2\) , therefore

\[
| \mu (g) | \leqslant | \mu (g _ {1}) | + | \mu (g _ {2}) | \leqslant L (f _ {1}) + L (f _ {2});
\]

since \(|\mu(g)|\) is arbitrarily close to \(L(f_1 + f_2)\) , we have

\[
L (f _ {1} + f _ {2}) \leqslant L (f _ {1}) + L (f _ {2}),
\]

which completes the proof of our assertion.

When \(\mu\) is a real measure, it follows from formula (9) that \(|\mu| \leqslant L\) ; on the other hand, by virtue of the last part of No.~5, if \(g \in \mathcal{K}(\mathbf{X}; \mathbf{C})\) and \(|g| \leqslant f \in \mathcal{K}_{+}(\mathbf{X})\) then \(|\int g d\mu| \leqslant \int |g| \cdot d|\mu| \leqslant \int f d|\mu|\) , therefore by definition \(L \leqslant |\mu|\) , in other words \(L = |\mu|\) .

We denote again by \(|\mu|\) the positive measure L for any complex measure \(\mu\) , and we say that \(|\mu|\) is the absolute value of \(\mu\) . The definition of \(|\mu|\) can therefore be written

\[
| \mu | (f) = \sup _ {| g | \leqslant f, g \in \mathcal{K} (\mathrm{X}; \mathbf {C})} | \mu (g) |,\tag{12}
\]

consequently, for every function \(g \in \mathcal{K}(\mathrm{X};\mathbf{C})\)

\[
\left| \int g d \mu \right| \leqslant \int | g | d | \mu |.\tag{13}
\]

It is clear that for every scalar \(\alpha \in \mathbf{C}\) and every measure \(\mu\) on \(X\) ,

\[
| \alpha \mu | = | \alpha | \cdot | \mu |.\tag{14}
\]

On the other hand, if \(\mu\) and \(\nu\) are two measures on \(X\) , \(f\) is a function in \(\mathcal{K}_{+}(X)\) , and \(g\) is a function in \(\mathcal{K}(X; \mathbf{C})\) such that \(|g| \leqslant f\) , then

\[
\left| \int g d (\mu + \nu) \right| = \left| \int g d \mu + \int g d \nu \right| \leqslant \int f d | \mu | + \int f d | \nu |,
\]

whence

\[
| \mu + \nu | \leqslant | \mu | + | \nu |.\tag{15}
\]

With the same notations, the relations \(|g| \leqslant f\) and \(|\overline{g}| \leqslant f\) are equivalent, therefore

\[
| \overline {{{\mu}}} | = | \mu |.\tag{16}
\]

It follows from (14), (15) and (16) that

\[
| \mathcal {R} \mu | \leqslant | \mu |, \quad | \mathcal {I} \mu | \leqslant | \mu |, \quad | \mu | \leqslant | \mathcal {R} \mu | + | \mathcal {I} \mu |.\tag{17}
\]

PROPOSITION 8. --- If \(\mu\) is a measure on X then, for every function \(h \in \mathcal{C}(\mathrm{X};\mathbf{C})\) ,

\[
\left| h \cdot \mu \right| \leqslant \left| h \right| \cdot \left| \mu \right|.\tag{18}
\]

For, if \(f \in \mathcal{K}_+(\mathbf{X})\) and if \(g \in \mathcal{K}(\mathbf{X};\mathbf{C})\) is such that \(|g| \leqslant f\) , then, by (13), \(|\int gh d\mu| \leqslant \int |gh| d|\mu| \leqslant \int f|h| d|\mu|\) , which proves (18).

